
\documentclass{article}
\usepackage[utf8]{inputenc}
\usepackage{geometry}
\geometry{a4paper, margin=1in}
\usepackage{titling}

\title{MEMORANDUM}
\author{Team \#2607504}
\date{February 4, 2026}

\begin{document}

\section*{MEMORANDUM}
\noindent\textbf{To:} Executive Producers, \emph{Dancing with the Stars}\\
\textbf{From:} Team \#2607504\\
\textbf{Date:} February 4, 2026\\
\textbf{Subject:} Proposal for System Upgrade: Maximizing Drama \& Integrity with ``ACDW-B3''

\hrule height 1pt
\vspace{1em}

\subsection*{Executive Summary: The ``ACDW-B3'' Solution}
We recommend transitioning from the current rigid scoring rules to a dynamic system:
\textbf{ACDW-B3 (Adaptive Concave Diminishing-returns Weighted Bottom-3)}.
Our analysis of 34 seasons reveals that current rules are unstable (agreeing only 38.2\% of the time) and prone to ``fan hijackings.'' The proposed system achieves a \textbf{90.5\% comprehensive fairness score}, significantly outperforming the current best options by balancing technical merit with audience engagement.

\vspace{-0.5em}

\subsection*{The Problem: Why Change Now?}
\noindent\textbf{The ``Superstar'' Distortion:} Under the Percentage Rule, massive fanbases allow celebrities (e.g., Bobby Bones) to mathematically eliminate the need for dance skills. Our data shows this rule has a 0.945 ``Pro-Fan Bias''.\\
\textbf{The ``Boring'' Safe Zone:} The Rank Rule protects talent but often creates stagnant leaderboards, lacking the jeopardy that drives viewership.

\vspace{-0.5em}

\subsection*{Strategic Insight: Casting ROI (Addressing Problem C Req.)}
Our Triple-Track Mixed-Effects Model reveals critical insights for budget allocation:
\begin{itemize}
    \item \textbf{The ``Partner'' Myth:} Statistical modeling proves that \textbf{professional partners have negligible impact} on judicial scoring or fan votes.
    \item \textbf{Casting Recommendations:}
    \begin{itemize}
        \item \textbf{Actors} are the safest investment for high technical scores.
        \item \textbf{Politicians} generate high engagement (bias) but lower technical quality.
        \item \textbf{Strategy:} Allocated budget should prioritize high-profile celebrities rather than competing for costly ``star'' professional dancers.
    \end{itemize}
\end{itemize}

\vspace{-0.5em}

\subsection*{The Solution: How ACDW-B3 Works}
\begin{enumerate}
    \item \textbf{The ``Superstar Cap'' (Fan Utility Compression)}\\
    \textbf{Concept:} We apply a ``diminishing returns'' curve ($p=0.55$) to raw vote counts.\\
    \textbf{Effect:} The first 10,000 votes matter more than the 1,000,000th vote. This ensures popular stars stay safe but cannot generate an insurmountable lead that renders the judges irrelevant.
    
    \item \textbf{The ``Consensus Meter'' (Dynamic Weighting)}\\
    \textbf{Concept:} A dynamic algorithm measures the agreement between Judges and Fans weekly.\\
    \textbf{Mechanism:}
    \begin{itemize}
        \item \textbf{High Agreement:} Fan weight increases (Let the audience decide).
        \item \textbf{High Conflict:} Judge weight increases automatically (up to 80\%).
    \end{itemize}
    \textbf{TV Opportunity:} Visualize this as an on-screen \textbf{``Tension Bar''} before results, showing viewers how divided the house is.
    
    \item \textbf{The ``Danger Zone'' Expansion (Bottom-3 + Judge Save)}\\
    \textbf{Concept:} Expand the elimination zone from Bottom-2 to Bottom-3.\\
    \textbf{Effect:} This pulls more couples into the spotlight, increasing dramatic tension.\\
    \textbf{The Save:} Judges explicitly save the dancer with the highest technical score among the bottom three. This removes accusations of bias while protecting excellence.
\end{enumerate}

\vspace{-0.5em}

\subsection*{Projected Impact \& Implementation}
\noindent\textbf{Risk Mitigation:} The ``Consensus Meter'' provides transparency, justifying why judge influence scales up during controversial weeks.\\
\textbf{Performance:} Simulations show this system retains \textbf{97.2\% alignment} with professional standards (Judge Alignment) while still respecting fan rankings in 76.3\% of cases.\\
\textbf{Next Step:} We recommend a pilot run of ACDW-B3 shadow-scoring alongside the current season to validate these metrics before a full public rollout.

\end{document}
